\section{Conclusion}

What began as an experiment with GraphRAG applied to six immunology articles has evolved into a domain-agnostic framework validated across two fundamentally different knowledge domains: biomedical drug discovery and event contract management. The trajectory of that evolution---from similarity-based graph assembly to comprehension-based extraction, and from a single-domain tool to a pluggable framework---is itself the central finding of this work.

Similarity-based approaches to knowledge graph construction, including GraphRAG and related embedding-retrieval methods, produce graphs that reflect lexical co-occurrence patterns in source text. These graphs are useful for information retrieval but do not capture the typed, directional, evidence-backed relationships that domain experts require. Epistract demonstrates that parallel LLM extraction agents, constrained by a domain schema and validated by domain-specific tools, can produce knowledge graphs grounded in meaning rather than statistical proximity. Across six drug discovery scenarios---neurogenetics, oncology, rare disease, immuno-oncology, cardiovascular inflammation, and GLP-1 competitive intelligence---the system achieved 100\% entity type coverage and 93\% relation type coverage with zero retraining. Applied to a corpus of 62 event planning contracts, the same pipeline infrastructure extracted 341 nodes and 663 edges, detecting 53 cross-contract conflicts---requiring only a new domain configuration package, with no modifications to the core pipeline.

The contribution is threefold. First, the method---agentic extraction with typed domain schemas, parallel agent dispatch, and a two-layer knowledge graph architecture (brute facts grounded to domain ontologies plus an epistemic super-domain that surfaces contradictions, gaps, and risks)---establishes a reusable pattern for building knowledge graphs from any document corpus. Second, the framework---with its three-layer architecture separating core pipeline, pluggable domain configurations, and example consumers---makes that method accessible to domain experts who can create new domains without modifying pipeline code. Third, the tool---an open-source, MIT-licensed Claude Code plugin installable in one command---makes the framework immediately usable, with pre-built domains for drug discovery and event contract management.

The human-AI collaboration model through which Epistract was built is a fourth, less conventional contribution. Production-quality scientific software---including the extraction pipeline, validated test scenarios across two domains, project documentation, and this paper---was constructed through intensive interactive sessions where human scientific judgment directed AI implementation capabilities. This is not a claim about AI autonomy; it is a claim about the productivity of a specific collaboration pattern that other teams can adopt regardless of domain.

Several directions remain open for future work. A domain registry would enable community-contributed domain packages, expanding the framework's coverage through collective effort rather than centralized development. Export to Neo4j or similar graph databases would enable production-scale queries and integration with existing knowledge management infrastructure. Combining the extracted knowledge graph with vector embeddings could support hybrid retrieval that leverages both structural relationships and semantic similarity. Finally, applying the framework to additional domains---regulatory compliance, legal contract analysis, technical documentation---would further test the generalizability of the pluggable architecture and expand the ecosystem of available domain configurations.
